\documentclass[../main.tex]{subfiles}
\graphicspath{{\subfix{../images/}}}
\begin{document}

\section{Appendix to Chapter \ref{chap:mbvqe}}

The following appendix is taken from \cite{schroeder2024deterministic}.

\subsection{Overview of other previous works}~\label{app:overview}

The work by Ferguson et al. \cite{ferguson2021measurement} inspired further investigations of the performance and resource requirements of other VQE instances in the MBQC framework. Proietti et al.~\cite{proietti2022native} studied an MBQC native way to implement the Quantum Approximate Optimization Algorithm (QAOA) \cite{farhi2014a}, which makes use of the fact that operations that are diagonal in the computational basis are exceptionally efficient in MBQC as ``input'' and  ``output'' qubits coincide \cite{browne2006one}. They tested the MBQC algorithm by simulating this method in the context of the k-max cut problem.
Qin et al. \cite{qin2023applicability} showed an MBQC version of the Hamiltonian variational ansatz \cite{wecker2015progress} that makes use of efficient 2-body operators implemented in the same sense as the QAOA operator. Combined with local rotations on the input/output qubits, it is possible to generate all Pauli rotations.
Chan et al. \cite{chan2023hybrid} included these efficient MBQC-based multi-qubit Pauli-rotations, also referred to as Pauli-gadgets, within a gate-based circuit. This work included experiments on the IBM 7-qubit machine that indicate a possible advantage of operating gate-based hardware in an MBQC manner.
Marqversen et al. \cite{marqversen2023applications} discuss and exhibit how MBQC can reduce hardware resources for applications such as VQE, demonstrate specific use cases, and introduce efficient ways to simulate MBQCs with tensor networks.
While Majumder et al. \cite{majumder2023variational} used almost the same MBQC pattern as in our work to generate their ansatz and acknowledged the importance of byproducts in MBVQE, but instead of correcting them, they utilized this randomness to learn probability distributions better in the context of generative modeling. A part of these works used post-selection to guarantee the convergence of the variational algorithm~\cite{ferguson2021measurement, proietti2022native}, while others used different measurement bases than $R(\theta)$ and achieved determinism~\cite{chan2023hybrid, huang2022near, marqversen2023applications}. While sufficient attention was not given to determinism in the initial works of MBVQE, recent papers (including ours) discuss the problem extensively for distinct settings and applications. 


\subsection{Graph state modifications}\label{appendix:graph}
Here, we recall some valuable relations of graph states that are used in this work. Local  complementations (LC) implements \cite{adcock2020mapping,van2004graphical} the local unitary transformation of the following form:
\begin{equation}
 \ket{G_{LC}} = \sqrt{-iX_i} \bigotimes_{j \in \mathcal{N}(i)} \sqrt{iZ_j} \ket{G}.
\end{equation}

\begin{figure}
    \centering
    \includegraphics[width=0.8\linewidth]{mbvqe/graphmod.pdf}{
    \phantomsubcaption\label{fig:lclu}
    \phantomsubcaption\label{fig:measpattern}
    }
    \caption{\textbf{Left}: Graph states linked through local complementation are local Clifford equivalent. In our experiments, we choose to implement the left graph since it can be directly executed on the IBM device. \textbf{Right}: By preparing the heavy-hex graph state and measuring the blue qubits in the illustrated bases, it is possible to generate a cluster state without the need for SWAP gates. However, it's important to note that the measurement outcomes of Pauli measurements on graph states are not graph states but rather local unitary equivalents. Despite this, these ``byproduct'' operators can be efficiently computed from the measurement outcomes of the auxiliary qubits and accounted for in subsequent calculations due to Gottesman–Knill theorem. This figure was taken from Ref.~\cite{schroeder2024deterministic}.}
    \label{fig:lc}
\end{figure}


If we want to implement a specific graph state, we could choose to implement the representative from the LC equivalence class with the lowest edge count or the representative that best suits the hardware topology.
That is shown in Fig.~\ref{fig:lclu} -- the edge between the two qubits in the middle layer requires SWAP gates if implemented naively on quantum hardware with heavy-hex connectivity. However, it is easy to see that this edge can be eliminated by local complementation on the parent-qubit. This graph can be directly mapped onto the respective hardware.

To generate a cluster state with constant gate depth from hardware featuring heavy-hex connectivity, we start by creating a heavy-hex graph state. 

It is known that graph states stay graph states up to specific local unitaries under Pauli measurements, and efficient graphical rules were developed to predict the post-measurement state \cite{hein2006entanglement, anders2006fast}. 
By using these rules and minding the commutation relation of measurement projectors and the emerging local unitaries (see Table \ref{tab:commutation_relations}), it is possible to design an effective measurement pattern illustrated in Fig.~\ref{fig:measpattern}. 

With this scheme, each \textit{honeycomb} becomes a square of the cluster state. Hence, on \texttt{ibmq\_ehingen}, we can generate a $2\times3$ cluster state, and on the 127-qubit machine \texttt{ibmq\_nazca} $4\times7$.

\begin{table}
    \centering
    \begin{tabular}{ccc}
         $P_{x,\pm} \sigma_z = \sigma_z P_{x,\mp}$, & \quad &  $P_{y,\pm} \sqrt{-i\sigma_z} = \sqrt{-i\sigma_z} P_{x,\pm}$, \\ 
         $P_{y,\pm} \sigma_z = \sigma_z P_{y,\mp}$, & \quad &  $P_{y,\pm} \sqrt{i\sigma_y} = \sqrt{i\sigma_y} P_{y,\pm}$, \\
         $P_{z,\pm} \sigma_z = \sigma_z P_{z,\pm}$, & \quad &  $P_{y,\pm} \sqrt{-i\sigma_y} = \sqrt{-i\sigma_y} P_{y,\pm}$, \\
         \quad & \quad &                                        $P_{y,\pm} \sqrt{i\sigma_z} = \sqrt{i\sigma_z}P_{x,\mp}$, \\
         
         \quad & \quad & \quad \\
         
         $P_{x,\pm} \sqrt{-i\sigma_z} = \sqrt{-i\sigma_z} P_{y,\mp}$,                & \quad &  $P_{z,\pm} \sqrt{-i\sigma_z} = \sqrt{-i\sigma_z} P_{z,\pm}$, \\
         $P_{x,\pm} \sqrt{i\sigma_y} = \sqrt{i\sigma_y} P_{z,\textcolor{red}{\mp}}$, & \quad &  $P_{z,\pm} \sqrt{i\sigma_y} = \sqrt{i\sigma_y} P_{x,\pm}$, \\
         $P_{x,\pm} \sqrt{-i\sigma_y} = \sqrt{-i\sigma_y} P_{z,\mp}$,                & \quad &  $P_{z,\pm} \sqrt{-i\sigma_y} = \sqrt{-i\sigma_y} P_{x,\textcolor{red}{\mp}}$, \\
         $P_{x,\pm} \sqrt{i\sigma_z} =\sqrt{i\sigma_z} P_{y,\mp}$,                   & \quad &  $P_{z,\pm} \sqrt{i\sigma_z} = \sqrt{i\sigma_z} P_{z,\pm}$. \\
    \end{tabular}
    \caption{The commutation relations play an essential role in simulating post-measurement states of multiple Pauli measurements on a graph state, as the resulting state is a graph state up to local Clifford operations. These additional operations alter successive measurements but in an efficiently simulable manner. The precise form of these operators and the table can be found in \cite{hein2006entanglement}. We opted to include the commutation relations in the appendix due to the presence of typos in the highlighted equations of the original source. }
    \label{tab:commutation_relations}
\end{table}




\subsection{Clifford-data regression}\label{appendix:cdr}
Due to noise in the quantum hardware, expectation values extracted from circuit measurements have to be post-processed.
Several error-mitigation strategies exist, such as zero-noise-extrapolation~\cite{PhysRevLett.119.180509}, self-mitigation~\cite{ARahman:2022tkr}, or probabilistic error cancellation~\cite{van2023probabilistic,gupta2023probabilistic}.
For our experiments, we choose Clifford-data regression~\cite{Czarnik2021errormitigation} since dynamic circuits can be easily incorporated into this framework.

Given a parametrized circuit $\mathcal{C}(\theta_1,\dotsc \theta_n)$ and a Hamiltonian $H$, of which we try to estimate the expectation value $\bra{\psi}H\ket{\psi}$ at some given point $(\theta_1,\dotsc \theta_n)$, we first evaluate the circuit at near-Clifford points, i.e., we substitute some of the parameters by multiples of $\pi/2$, $\theta_i = n\cdot\pi/2,\ n\in \mathbb{N}$, for which the circuit can be simulated efficiently classically.
At these so-called regression points, we can find a relationship between the measured, noisy expectation values $\expval{H}{\psi}_\mathrm{noisy}$, and the exact, simulated ones,
\begin{equation}
    \expval{H}{\psi}_\mathrm{exact}=f(\expval{H}{\psi}_\mathrm{ noisy},a_0,\dotsc a_n),
\end{equation}
where $a_i$ are model parameters.
In this work, we assume a linear relationship, i.e.,
\begin{equation}
    \expval{H}{\psi}_\mathrm{exact}=a_0 \expval{H}{\psi}_\mathrm{noisy}+a_1,
\end{equation}
which can be justified by the assumption of the existence of a global noise channel \cite{Czarnik2021errormitigation}.


\section{Appendix to Chapter \ref{chap:chem}}

The following appendix is taken from \cite{amsler2023classical} and \cite{kiser2024classical}.

\subsection{Molecular strctures}


For the calculations on singlet and triplet molecular oxygen, and ozone presented in Sec.~\ref{sec:relative_energies}, Table~\ref{tab:results_exp_geo} and Table~\ref{tab:results_dft_geo}, we use experimental geometries from \cite{tanaka1970coriolis, krupenie1972spectrum}:
\begin{table}[h!]
\ce{^3O2}: \begin{tabular}{|c| S[table-format=+1.7] | S[table-format=+1.7] |S[table-format=+1.7]|}\hline
{} & {$x/\AA{}$} & {$y/\AA{}$} & {$z/\AA{}$} \\ \hline
\ce{O} & 0.0 & 0.0 &  0.603760\\
\ce{O} & 0.0 & 0.0 & -0.603760\\ \hline
\end{tabular} \nonumber \\
\ce{^1O2}: \begin{tabular}{|c| S[table-format=+1.7] | S[table-format=+1.7] |S[table-format=+1.7]|}\hline
\ce{O} & 0.0 & 0.0 &  0.607800\\
\ce{O} & 0.0 & 0.0 & -0.607800\\ \hline
\end{tabular}  \\
\ce{^{\phantom{3}}O3}: \begin{tabular}{|c| S[table-format=+1.7] | S[table-format=+1.7] |S[table-format=+1.7]|}\hline
\ce{O} & 0.0 & 0.0 &  0.0\\
\ce{O} & 0.0 & 0.0 &  1.2717000\\ 
\ce{O} & 1.1383850 & 0.0 & 1.8385340\\ \hline
\end{tabular}\nonumber 
\caption{Table was taken from \cite{amsler2023classical}.}\label{eq:exp_geo} 
\end{table}

We use geometries optimized at the density functional theory (DFT) level (UB3LYP/def2-QZVPP) as well, resulting in the xyz-structures:  
\begin{table}[h!]
\ce{^3O2}: \begin{tabular}{|c| S[table-format=+1.7] | S[table-format=+1.7] |S[table-format=+1.7]|}\hline
{} & {$x/\AA{}$} & {$y/\AA{}$} & {$z/\AA{}$} \\ \hline
\ce{O} & 0.0 & 0.0 &  0.6020120\\
\ce{O} & 0.0 & 0.0 & -0.6020120\\ \hline
\end{tabular} \nonumber \\
\ce{^1O2}: \begin{tabular}{|c| S[table-format=+1.7] | S[table-format=+1.7] |S[table-format=+1.7]|}\hline
\ce{O} & 0.0 & 0.0 &  0.6021008\\
\ce{O} & 0.0 & 0.0 & -0.6021008\\ \hline
\end{tabular}  \\
\ce{^{\phantom{3}}O3}: \begin{tabular}{|c| S[table-format=+1.7] | S[table-format=+1.7] |S[table-format=+1.7]|}\hline
\ce{O} &  3.1946229 & 3.9056962 &  0.0\\
\ce{O} & -2.1046353 & 3.2467339 &  0.0\\ 
\ce{O} & -1.0278818 & 3.9270800 & 0.0\\ \hline
\end{tabular}\nonumber
\caption{Table was taken from \cite{amsler2023classical}.}\label{eq:dft_geo}
\end{table}




\begin{table}[h!]
\begin{tabular}{c|c|c|c|c|c|c} 
  &  & HF & MP2 & BP86 & B3LYP & M06-2X \\ \hline
 \ce{^3O2} & reference  & UHF & UHF & UKS & UKS & UKS \\
 & $r_\mathrm{\ce{O}-\ce{O}}$ [pm] & 115.79 & 121.90 & 122.03 & 120.43 & 118.79 \\
 & ZPVE [kJ/mol] & 11.78 & 8.84 & 9.24 & 9.78 & 10.55 \\ \hline
  \ce{^1O2} & reference  & UHF & RHF & UKS & UKS & UKS \\
 & $r_\mathrm{\ce{O}-\ce{O}}$ [pm] & 115.61 & 124.34 & 122.12 & 120.46 & 118.73 \\ 
 & ZPVE [kJ/mol] & 11.80 & 7.72 & 9.20 & 9.75 & 10.54 \\ \hline
\ce{O3} &  reference  & UHF & RHF & RKS (=UKS) & UKS & UKS \\
 & $r_\mathrm{\ce{O}-\ce{O}}$ [pm] & 127.54 & 127.81 & 127.55 & 127.41 & 125.69 \\
 & $\theta_{\ce{O}-\ce{O}-\ce{O}}$ [$^{\circ}$] & 118.14 & 116.76 & 118.14 & 116.55 & 116.05 \\
 & ZPVE [kJ/mol] & 17.66 & 24.69 & 17.63 & 16.87 & 18.57 \\ 
\end{tabular} 
\caption{A benchmark on the method dependence of structural parameters as well as zero-point vibrational energies (ZPVE). All calculations are performed using the cc-pVQZ basis. This Table was taken from \cite{amsler2023classical}. \label{dft-struct-zpe-comparison}}
 \end{table}

\begin{table}[!]
\begin{tabular}{l|l|l|l}                  & \ce{H4}                    & \ce{O2} \& \ce{O3}              & Hubbard system (\ce{CuBr2}) \\ \hline
Number of walkers $N_\mathrm{W}$            & $10^3$ & $10^3$ & $10^3$                       \\ \hline
Number of blocks  $N_\mathrm{B}$             & $10^3$ & $10^4$ & $10^4$                \\ \hline
Time steps per block           & 10                    & 10                    & 10                     \\ \hline
Equilibration time (in blocks) & 100                   & 1000                  & 3000                   \\ \hline
Time step $\Delta\tau$ (a.u.)               & $0.005$                 & $0.005$                 & $0.005$                  \\ 
\end{tabular}
\caption{Parameters used in the AFQMC calculations. Parameters not reported here are set to the default setting of ipie \cite{malone2022a}. This table was taken from \cite{amsler2023classical}. \label{tab:comp_details}}
\end{table}
 
\begin{table}[h!]
\centering
\resizebox{\textwidth}{!}{
\begin{tabular}{p{8mm}p{30mm}p{25mm}p{25mm}p{25mm}p{20mm}p{20mm}p{0mm}p{0mm}}
\hline
Basis & Method &  $E(\ce{^3O_2})$ [Ha] & $E(\ce{^1O_2})$ [Ha] & $E(\ce{O_3})$ [Ha] & $\Delta E (\ce{^1O_2})$ [kJ/mol] & $\Delta E (\ce{O_3})$ [kJ/mol]       &  &   \\
\hline\hline
& Experiment \cite{krupenie1972spectrum, weast1983} & - & - & - & 94.7 & 143.2 \\
\hline\hline
\multirow{6}{*}{\rotatebox{90}{\usebox0}}
& R(O)HF & -149.6083 & -149.5414 & -224.2657 & 175.7 & 385.3 & &  \\ \cline{2-9}
& CCSD(T) & -149.9892 & -149.9405 & -224.9149 & 127.9 & 180.9 & &  \\ \cline{2-9}
& CASSCF & -149.7087 & -149.6757 & -224.4976 & 86.7 & 171.9 & &  \\ \cline{2-9}
& NEVPT2 & -149.9579 & -149.9209 & -224.8759 & 96.9 & 159.8 & & \\ \cline{2-9}
& AFQMC(HF) & -149.9827(7) & -149.9403(6) & -224.9116(14) & 111.2(25) & 164.0(41) & &  \\ \cline{2-9}
& AFQMC(CASSCF) & -149.9912(3) & -149.9500(2) & -224.9188(4) & 108.0(10) & 178.3(13) & &  \\ \hline
\hline

\multirow{6}{*}{\rotatebox{90}{\usebox1}}
& R(O)HF & -149.6528 & -149.5877 & -224.3405 & 171.1 & 364.4 & &  \\ \cline{2-9}
& CCSD(T) & -150.1536 & -150.1058 & -225.1700 & 125.4 & 158.5 & &  \\ \cline{2-9}
& CASSCF & -149.7519 & -149.7192 & -224.5688 & 86.1 & 155.2 & &  \\ \cline{2-9}
& NEVPT2 & -150.1138 & -150.0782 & -225.1179 & 93.5 & 138.8 & & \\ \cline{2-9}
& AFQMC(HF) & -150.1494(7) & -150.1055(8) & -225.1651(10) & 115.3(28) & 155.1(31) & &  \\ \cline{2-9}
& AFQMC(CAS) & -150.1554(4) & -150.1146(3) & -225.1740(4) & 106.9(13) & 155.1(15) & &  \\ \hline
\hline
\multirow{8}{*}{\rotatebox{90}{\usebox2}}
& R(O)HF & -149.6643 & -149.5995 & -224.3585 & 170.0 & 362.1 & &  \\ \cline{2-9}
& UPBE & -150.2569 & -150.2431 & -225.3366 & 36.3 & 128.0 & &  \\ \cline{2-9}
& UB3LYP & -150.3376 & -150.3214 & -225.4366 & 42.7 & 183.2 & &  \\ \cline{2-9}
& CCSD(T) & -150.2343 & -150.1871 & -225.2929  & 124.0 & 153.8 & &  \\ \cline{2-9}
& CASSCF & -149.7633 & -149.7306 & -224.5866 & 85.9 & 153.2  & &  \\ \cline{2-9}
& NEVPT2 & -150.1928 & -150.1577 & -224.2378 & 92.0 & 134.8 & & \\ \cline{2-9}
& AFQMC(HF) & -150.2279(7) & -150.1880(7) & -225.2901(12) & 104.7(27) & 135.8(36) & &  \\ \cline{2-9}
& AFQMC(CASSCF) & -150.2354(4)  & -150.1962(3) & -225.2966(4) & 102.9(13) & 148.4(14) & &  \\ \hline
\hline

\multirow{4}{*}{\rotatebox{90}{\usebox3}}
& CCSD(T) & -150.2879 & -150.2409 & -225.3742 &123.4 &151.3 & &  \\ \cline{2-9}
& CASSCF & -149.7662 & -149.7334 & -224.5911 &86.3 & 152.8 & &  \\ \cline{2-9}
& NEVPT2 & -150.2450 & -150.2102 & -225.3169 & 91.4 & 132.9 & & \\ \cline{2-9}
& AFQMC(CASSCF) & -152.2884(7) & -150.2502(6) & -225.3776(8) &100.4(24) & 144.4(27) & &  \\ \hline
\hline
\end{tabular}
}
\caption{Results using the experimental geometries given in (\ref{eq:exp_geo}) for various methods and basis sets together with the extrapolated results to the complete basis set (CBS) limit. For all AFQMC(CASSCF), NEVPT2, and CASSCF calculations, (8e, 6o) and (12e, 9o) active spaces were used for singlet and triplet molecular oxygen, and ozone, respectively. All other calculations were done in the canonical MO basis (R(O)HF). For ozone, the experimental value was corrected by the ZPVEs as described in the main text. Tables were taken from \cite{amsler2023classical}.}
\label{tab:results_exp_geo}
\end{table}



\begin{table}[h!]
\centering 
\resizebox{\textwidth}{!}{
\begin{tabular}{p{30mm}p{25mm}p{25mm}p{25mm}p{20mm}p{20mm}p{0mm}p{0mm}}
\hline
Method &  $E(\ce{^3O_2})$ [Ha] & $E(\ce{^1O_2})$ [Ha] & $E(\ce{O_3})$ [Ha] & $\Delta E (\ce{^1O_2})$ [kJ/mol] & $\Delta E (\ce{O_3})$ [kJ/mol]       &  &   \\
\hline\hline
Experiment \cite{krupenie1972spectrum, weast1983} & - & - & - & 94.7 & 143.2 \\
\hline\hline
R(O)HF & -149.6649 & -149.6018 & -224.3578 & 165.6  & 366.4 & &  \\ \cline{1-8}
UPBE & -150.2568  & -150.2428 & -225.3366 & 36.9 & 127.6 & &  \\ \cline{1-8}
UB3LYP & -150.3377 & -150.3215 & -225.4366 & 42.4 & 183.5 & &  \\ \cline{1-8}
CCSD(T) & -150.2344 & -150.1869 & -225.2929 & 124.7 & 154.0 & &  \\ \cline{1-8}
CASSCF & -149.7632 & -149.7301 & -224.5867 & 87.0 & 152.8 & &  \\ \cline{1-8}
NEVPT2 & -150.1928 & -150.1576 & -225.2378 & 92.4 & 134.7 & & \\ \cline{1-8}
AFQMC(CASSCF)& -150.2356(4)  & -150.1957(3) & -225.2974(4) & 104.9(12) & 147.1(14) & &  \\ \hline
\hline

\end{tabular}
}
\caption{Results using the DFT-optimized geometries given in (\ref{eq:dft_geo}) for various methods using the def2-QZVPP basis set for DFT (B3LYP and PBE) and cc-pVQZ otherwise. For all AFQMC(CASSCF), NEVPT2, and CASSCF calculations, (8e, 6o) and (12e, 9o) active spaces were used for singlet and triplet molecular oxygen, and ozone, respectively. All other calculations were done in the canonical MO basis (R(O)HF). For ozone, the experimental value was corrected by the ZPVEs as described in the main text. Table was taken from \cite{amsler2023classical}.}
\label{tab:results_dft_geo}
\end{table}



\section{Appendix to Chapter \ref{chap:qnlp}}
\subsection{Hyperparameters}\label{sec:apn-hyp}

The following hyperparameters were taken from Ref.~\cite{tiblias2025an}.

A randomized search is performed to find the best hyperparameters for each model and data set. $50$ configurations, train them on the task, and evaluate their performance on the development set. To avoid overfitting, we perform early stopping on the training if the development loss does not decrease for five consecutive epochs. Throughout all experiments, we utilize the \textit{Adam} optimizer provided by PyTorch. Since our largest model, HAM, has at most $\sim500k$ parameters (dictated by the embedding space), we limit the random search to configurations with fewer than or equal to the number of parameters to ensure a fair comparison. In practice, we find RNNs, LSTMs, and CNNs perform very well with as few as $40k$ parameters. What follows is a list of all hyperparameters we evaluated:
\begin{itemize}
    \item \textbf{Batch size:} $[64, 128, 256]$;
    \item \textbf{Learning rate:} $[10^{-2}, 10^{-3}, 10^{-4}]$; 
    \item \textbf{Hidden size:} For RNNs and LSTMs, the size of the hidden representation $[100, 300, 500]$;
    \item \textbf{Layers:} For RNNs and LSTMs, the number of layers of the recurrent block $[1, 4, 8]$. For MLP, the total number of layers including input and output $[ 3, 4, 5 ]$. For CNNs, the total number of convolutional layers $[ 3, 4, 5 ]$. For HAM, PEFF, SIM and CIRC, the number of repeated applications of the ansatz, analogous to the number of layers $[8, 16, 32]$;
    \item \textbf{Channels:} For CNNs, the number of output channels of each convolutional layer $[ 8, 16, 32, 64, 128]$;
    \item \textbf{Kernel size:} For CNNs, the kernel dimension of each convolutional layer $[ 3, 5, 7 ]$;
    \item \textbf{Circuit:} The ansatz $U_\theta$ used to prepare state $\psi$. We experiment with three hardware-efficient ans\"atze inspired by circuits from \cite{Sim2019circ} and illustrated in Figure \ref{fig:ansz}. These ans\"atze are designed to explore different levels of entanglement (the quantum counterpart of correlation) and gate complexity. The first is a non-entangling circuit composed of single-qubit rotations, exhibiting linear complexity with respect to the number of qubits. The second is an entangling circuit arranged in a ring configuration, also scaling linearly in gate count. The third is an all-to-all entangling circuit, which scales quadratically. $[\texttt{Baseline, Ring, All-to-all}]$.
    \item \textbf{\# Pauli strings:} For SIM, the number of Pauli strings composing the Hamiltonian $[50, 100, 500, 1000]$.
\end{itemize}
Tables \ref{tab:hyp-baseline} and \ref{tab:hyp-ham} show the final selection of best-performing hyperparameters used in the main experiment of Section \ref{sec:qnlpexpres}.

\begin{figure}[h!]
    \centering
    % Circ base
    \begin{subfloat}[] {
        \begin{tikzpicture}
        \node[scale=0.6]{
        \begin{quantikz}[column sep=0.3cm]
 &  & \gate{R_X}\gategroup[3,steps=3,style={dashed,rounded corners, inner xsep=2pt}, background]{{}}    & \gate{R_Y} & \gate{R_Z} & & \\
 &  & \gate{R_X}    & \gate{R_Y} & \gate{R_Z} & & \\
 &  & \gate{R_X}    & \gate{R_Y} & \gate{R_Z} & &
\end{quantikz} %\ctrl{1}
        };
        \end{tikzpicture}
        }
    \end{subfloat}
    % Circ 14
    \begin{subfloat}[] {
        \begin{tikzpicture}
        \node[scale=0.4]{
        \begin{quantikz}[column sep=0.3cm] %\gategroup[3,steps=5,style={dashed,rounded corners, inner xsep=2pt}, background]{{}}
 &  & \gate{R_Y}\gategroup[3,steps=8,style={dashed,rounded corners, inner xsep=2pt}, background]{{}}    & \ctrl{1}   &              & \gate{R_X}    & \gate{R_Y}    & \ctrl{1}   &              & \gate{R_X}    & & \\
 &  & \gate{R_Y}    & \gate{R_X} & \ctrl{1}     &               & \gate{R_Y}    & \gate{R_X} & \ctrl{1}     &               & & \\
 &  & \gate{R_Y}    &            & \gate{R_X}   & \ctrl{-2}     & \gate{R_Y}    &            & \gate{R_X}   & \ctrl{-2}     & &
\end{quantikz} %\ctrl{1}
        };
        \end{tikzpicture}
        }
    \end{subfloat}
    % Circ 6
    \begin{subfloat}[] {
        \begin{tikzpicture}
        \node[scale=0.4]{
        \begin{quantikz}[column sep=0.3cm]   
 &  & \gate{R_X}\gategroup[3,steps=10,style={dashed,rounded corners, inner xsep=2pt}, background]{{}}     & \gate{R_Z} & \ctrl{1}   & \ctrl{2}    & \gate{R_X}    &               & \gate{R_X}    &               & \gate{R_X}    & \gate{R_Z} & & \\
 &  & \gate{R_X}    & \gate{R_Z} & \gate{R_X} &             & \ctrl{-1}     & \ctrl{1}      &               & \gate{R_X}    & \gate{R_X}    & \gate{R_Z} & &\\
 &  & \gate{R_X}    & \gate{R_Z} &            & \gate{R_X}  &               & \gate{R_X}    & \ctrl{-2}     & \ctrl{-1}     & \gate{R_X}    & \gate{R_Z} & &
\end{quantikz} %\ctrl{1}
        };
        \end{tikzpicture}
        }
    \end{subfloat}
    \caption{
    Circuit ansätze explored in the experiment. For the sake of visualization, the image shows ansätze for $n=3$. Since the experiments use different values, we extend the patterns to act on more qubits. 
    (a) Non-entangling ansatz
    (b) Ring ansatz
    (c) All-to-all ansatz. (b) and (c) are respectively circuit 14 and 6 from~\cite{Sim2019circ}. This figure was taken from Ref.~\cite{tiblias2025an}.}
    \label{fig:ansz}
\end{figure}

\begin{table*}[h!]
\centering
\begin{tabular}{lccc}
\toprule
\multicolumn{3}{c}{\textbf{LOG}} \\
\textbf{Dataset} & \textbf{Batch size} & \textbf{Learning rate} \\
\midrule
SST2 & 64 & $10^{-3}$ \\
IMDb & 64 & $10^{-2}$ \\
AG News & 64 & $10^{-2}$ \\
MNIST2 & 128 & $10^{-3}$ \\
CIFAR2 & 256 & $10^{-3}$ \\
Fashion & 64 & $10^{-3}$ \\
\bottomrule
\end{tabular}
\begin{tabular}{lccccc}
\toprule
\multicolumn{5}{c}{\textbf{MLP}} \\
\textbf{Dataset} & \textbf{Batch size} & \textbf{Learning rate} & \textbf{Hidden size} & \textbf{Layers} \\
\midrule
SST2 & 64 & $10^{-4}$ & 300 & 3 \\
IMDb & 64 & $10^{-3}$ & 300 & 3 \\
AG News & 64 & $10^{-3}$ & 100 & 3 \\
MNIST2 & 256 & $10^{-4}$ & 100 & 3 \\
CIFAR2 & 128 & $10^{-4}$ & 100 & 3 \\
Fashion & 256 & $10^{-3}$ & 100 & 3 \\
\bottomrule
\end{tabular}
\begin{tabular}{lccccc}
\toprule
\multicolumn{5}{c}{\textbf{LSTM}} \\
\textbf{Dataset} & \textbf{Batch size} & \textbf{Learning rate} & \textbf{Hidden size} & \textbf{Layers} \\
\midrule
SST2 & 64 & $10^{-2}$ & 100 & 2 \\
IMDb & 128 & $10^{-2}$ & 100 & 2 \\
AG News & 128 & $10^{-3}$ & 100 & 2 \\
\bottomrule
\end{tabular}
\begin{tabular}{lccccc}
\toprule
\multicolumn{5}{c}{\textbf{RNN}} \\
\textbf{Dataset} & \textbf{Batch size} & \textbf{Learning rate} & \textbf{Hidden size} & \textbf{Layers} \\
\midrule
SST2 & 64 & $10^{-4}$ & 100 & 1 \\
IMDb & 256 & $10^{-4}$ & 100 & 2 \\
AG News & 256 & $10^{-4}$ & 100 & 1 \\
\bottomrule
\end{tabular}
\begin{tabular}{lcccccc}
\toprule
\multicolumn{6}{c}{\textbf{CNN}} \\
\textbf{Dataset} & \textbf{Batch size} & \textbf{Learning rate} & \textbf{Layers} & \textbf{Channels} & \textbf{Kernel Size} \\
\midrule
MNIST2 & 64 & $10^{-3}$ & 5 & 16 & 5 \\
CIFAR2 & 64 & $10^{-4}$ & 4 & 32 & 5 \\
Fashion & 64 & $10^{-3}$ & 4 & 16 & 3 \\
\bottomrule
\end{tabular}
\caption{Best hyperparameters across all baseline models and tasks}
\label{tab:hyp-baseline}
\end{table*}
\begin{table*}
\centering
\begin{tabular}{lccccc}
\toprule
\multicolumn{5}{c}{\textbf{CIRC}} \\
\textbf{Dataset} & \textbf{Batch size} & \textbf{Learning rate} & \textbf{Circuit} & \textbf{Layers} \\
\midrule
SST2 & 256 & $10^{-3}$ & \texttt{All-to-all} & 8 \\
AG News & 128 & $10^{-3}$ & \texttt{Baseline} & 16 \\
IMDb & 128 & $10^{-3}$ & \texttt{All-to-all} & 16 \\
MNIST2 & 128 & $10^{-3}$ & \texttt{Baseline} & 8 \\
CIFAR2 & 64 & $10^{-3}$ & \texttt{Baseline} & 8 \\
Fashion & 256 & $10^{-3}$ & \texttt{Ring} & 32 \\
\bottomrule
\end{tabular}
\begin{tabular}{lccccc}
\toprule
\multicolumn{5}{c}{\textbf{HAM}} \\
\textbf{Dataset} & \textbf{Batch size} & \textbf{Learning rate} & \textbf{Circuit} & \textbf{Layers} \\
\midrule
SST2 & 128 & $10^{-3}$ & \texttt{Ring} & 8 \\
IMDb & 256 & $10^{-3}$ & \texttt{All-to-all} & 32 \\
MNIST2 & 256 & $10^{-2}$ & \texttt{Baseline} & 32 \\
CIFAR2 & 64 & $10^{-3}$ & \texttt{Ring} & 8 \\
\bottomrule
\end{tabular}
\begin{tabular}{lccccc}
\toprule
\multicolumn{5}{c}{\textbf{PEFF}} \\
\textbf{Dataset} & \textbf{Batch size} & \textbf{Learning rate} & \textbf{Circuit} & \textbf{Layers} \\
\midrule
SST2 & 64 & $10^{-2}$ & \texttt{All-to-all} & 8 \\
IMDb & 64 & $10^{-2}$ & \texttt{All-to-all} & 8 \\
MNIST2 & 64 & $10^{-3}$ & \texttt{All-to-all} & 8 \\
CIFAR2 & 256 & $10^{-2}$ & \texttt{Baseline} & 16 \\
\bottomrule
\end{tabular}
\begin{tabular}{lcccccc}
\toprule
\multicolumn{6}{c}{\textbf{SIM}} \\
\textbf{Dataset} & \textbf{Batch size} & \textbf{Learning rate} & \textbf{Circuit} & \textbf{Layers} & \textbf{\# Pauli strings} \\
\midrule
SST2 & 256 & $10^{-3}$ & \texttt{Ring} & 32 & 1000 \\
IMDb & 256 & $10^{-2}$ & \texttt{All-to-all} & 16 & 1000 \\
AG News & 128 & $10^{-3}$ & \texttt{All-to-all} & 16 & 1000 \\
MNIST2 & 256 & $10^{-2}$ & \texttt{Ring} & 32 & 1000 \\
CIFAR2 & 64 & $10^{-3}$ & \texttt{All-to-all} & 32 & 1000 \\
Fashion & 256 & $10^{-2}$ & \texttt{All-to-all} & 8 & 1000 \\
\bottomrule
\end{tabular}
\caption{Best hyperparameters across all Hamiltonian models and tasks. This table was taken from Ref.~\cite{tiblias2025an}.}
\label{tab:hyp-ham}
\end{table*}

\section{Code availability}

Code and data to reproduce the results and plots of Chapter~\ref{chap:mbvqe} and Chapter~\ref{chap:automaton} are available in the author's GitHub repository~\cite{schroeder2026quantumthesis}.

\end{document}